\centerline{\bf \Huge 9 класс}

\problem{\textbf{1}} \textbf{Критерии:} \\
\textit{7 баллов} $-$ Любое полное верное решение с найденным $x=27$, с доказательством существования ещё одного другого действительного корня и его примерным нахождением $x=1,15$.\\
\textit{6 баллов} $-$ Любое полное верное решение с найденным $x=27$ и с доказательством существования ещё одного другого действительного корня.\\
\textit{3 балла} $-$ Найден $x=27$ с обоснованием или с проверкой.\\
\textit{1 балл} $-$ Найден $x=27$ без обоснований и без проверки.\\
\textit{0 баллов} $-$ Любое неверное решение или неверный ответ.\\
\textit{0 баллов} $-$ Решение отсутствует.\\
\textit{0 баллов} $-$ При обоснованных подозрениях на списывание участником из любого источника (Интернет, учебное пособие, ключи олимпиады и т.д.) член жюри имеет право указать на такие фрагменты с помощью восклицательных знаков на полях. \\

\problem{\textbf{2}} \textbf{Критерии суммируются:} \\
\textit{0-4 балла} $-$ Решение А.\\
\textit{0-1 балл} $-$ Составление текстового условия Б.\\
\textit{0-2 балла} $-$ Любое правильное решение Б.\\
\textit{0 баллов} $-$ Любой ответ без объяснений.\\
\textit{0 баллов} $-$ Решение отсутствует.\\
\textit{0 баллов} $-$ При обоснованных подозрениях на списывание участником из любого источника (Интернет, учебное пособие, ключи олимпиады и т.д.) член жюри имеет право указать на такие фрагменты с помощью восклицательных знаков на полях. \\



\problem{\textbf{3}} \textbf{Критерии суммируются:} \\
\textit{2 балла} $-$ $1600+96\times\dfrac{120}{36}=96\times$количество коров => количество коров равно 20.\\
\textit{2 балла} $-$ Доказано, что на лугу (без роста травы) 1600 дневных порций коров.\\
\textit{3 балла} $-$ Доказано, что за день трава вырастает на $\dfrac{120}{36}$ дневных порций коров.\\
\textit{0 баллов} $-$ Любой ответ без объяснений.\\
\textit{0 баллов} $-$ Решение неверное, продвижения отсутствуют.\\
\textit{0 баллов} $-$ Решение отсутствует.\\
\textit{0 баллов} $-$ При обоснованных подозрениях на списывание участником из любого источника (Интернет, учебное пособие, ключи олимпиады и т.д.) член жюри имеет право указать на такие фрагменты с помощью восклицательных знаков на полях. \\

\textbf{ИЛИ:}
\textit{7 баллов} $-$ Другое полное верное решение.\\



\problem{\textbf{4}} \textbf{Критерии суммируются:} \\
\textit{3-4 баллов} $-$ Не доказано, что отрезок, соединяющий середины оснований, проходит через точку пересечения диагоналей трапеции, но в дальнейшем решение верное.\\
\textit{3 балла} $-$ Доказано, что отрезок, соединяющий середины оснований, проходит через точку пересечения диагоналей трапеции.\\
\textit{0 баллов} $-$ Любой ответ без объяснений.\\
\textit{0 баллов} $-$ Решение неверное, продвижения отсутствуют.\\
\textit{0 баллов} $-$ Решение отсутствует.\\
\textit{0 баллов} $-$ При обоснованных подозрениях на списывание участником из любого источника (Интернет, учебное пособие, ключи олимпиады и т.д.) член жюри имеет право указать на такие фрагменты с помощью восклицательных знаков на полях.\\

\textbf{ИЛИ:}
\textit{7 баллов} $-$ Полное верное решение.\\


\problem{\textbf{5}} \textbf{Критерии:} \\
\textit{7 баллов} $-$ Полное верное решение.\\
\textit{0 баллов} $-$ Любой ответ без объяснений.\\
\textit{0 баллов} $-$ Решение неверное, продвижения отсутствуют.\\
\textit{0 баллов} $-$ Решение отсутствует.\\
\textit{0 баллов} $-$ При обоснованных подозрениях на списывание участником из любого источника (Интернет, учебное пособие, ключи олимпиады и т.д.) член жюри имеет право указать на такие фрагменты с помощью восклицательных знаков на полях. \\